\chapter{Indian Sociologists on Caste: Summary and Critique}

\section{Summarizing the Four Scholars}

\begin{KeyConcept}
\begin{itemize}
\item \textbf{G. S. Ghurye} --- follows the \textbf{Orientalist} line: his six features of caste were the discovery of the Orientalists (who first translated the Sanskrit texts); his views were a \textbf{reaction} to the Orientalist idea of caste as responsible for 'underdevelopment'. He held caste \textbf{integrates} us (part of our culture). His account is \textbf{Indological, Evolutionary} (evolutionism from Darwin/Spencer) \textbf{and Comparative}.
\item \textbf{M. N. Srinivas} --- marks a \textbf{clear departure from Orientalism}: the first to start the \textbf{field view} officially/professionally/academically (earlier field workers like S. C. Roy existed, but not 'officially'). He used the \textbf{structural-functional approach} (whose origin is evolutionism, via Herbert Spencer), studying the \textbf{Coorgs} to show how joint family, marriage, rituals and caste \textbf{integrate} society. He developed \textbf{two concepts from field work: Sanskritization} (book/Varna view) \textbf{and dominant caste} (field/Jati view).
\item \textbf{Louis Dumont} --- an \textbf{Indologist} who believes \textbf{Brahmins are superior}: his line '\textbf{power is inferior to status}'. His two hierarchies are \textbf{pure/impure} and \textbf{status over power} (unlike Srinivas's ritual/secular). Against Weber (class, status, power independent) and Marx (power from money), Dumont says the top of the hierarchy has status, class and power together.
\item \textbf{Andre Beteille} --- 'slaps everyone': by showing \textbf{Brahmins are not homogeneous} (Smartha/Sri Vaishnava, further sub-divided), he exposed the \textbf{inadequacy of a sociology of caste running on Brahmanical lines} --- how can Brahmins be superior if they are themselves divided?
\end{itemize}
\end{KeyConcept}

\begin{CriticalPoint}
\begin{itemize}
\item \textbf{The common thread}: Ghurye, Srinivas and Dumont all end up \textbf{claiming Brahminical superiority} --- Ghurye (hierarchy with Brahmins at top; Vedic Aryan culture diffused south via Sage Agastya; Kartik replaced by Ganesh in the North), Srinivas (Sanskritization implies everyone wants to become Brahmin), Dumont (status belongs to Brahmins). Beteille is the critique of all three.
\item Summary keywords: Ghurye = Indological/Evolutionary/Comparative; Srinivas = empirical/fieldwork/structural-functional; Dumont = Indology + Structuralism (binaries: pure/impure, status/power); Beteille = the critic.
\end{itemize}
\end{CriticalPoint}

\section{Critique of Srinivas: Sanskritization}

\begin{CriticalPoint}
\begin{itemize}
\item \textbf{1. Sanskritization legitimises Brahminical supremacy even in the field view} --- it portrays Brahmins as the symbol/icon of civilisation, as if one cannot live without imbibing Brahminical culture. (Analogy: \textbf{Jevons' paradox} --- promoting fuel-efficient cars ends up burning more coal.)
\item \textbf{2. It ignores other religious traditions} --- Islam, Christianity and non-Hindu traditions. \textbf{Punjab got Islamized, not Sanskritized} (Gurudwara architecture is Indo-Islamic; Kabir's couplets are borrowed from the Quran; the ghazal mixes Sufi and Punjabi folk). It thus equates Indian culture with Hindu culture. \textbf{Yogendra Singh} talked of Islamization; the concept of '\textbf{multiple traditions}' was given by \textbf{S. C. Dube} (whom Yogendra Singh read). Sanskritization presents India as a \textbf{monolithic Hindu India}.
\item \textbf{3. The reference groups have changed} --- no longer upper-caste Hindu groups, but \textbf{secular, westernized groups} (an engineer at Facebook/Silicon Valley). So Sanskritization is \textbf{not empirically valid today}. Instead we see \textbf{Dalit Assertion} (Bhim Sena, Dalit Panther) --- asserting their own identity, not claiming Brahminical status. Beteille: in the light of \textbf{affirmative action and constitutional safeguards}, Dalits demand status \textbf{on humanitarianism, not at par with Brahmins}.
\item \textbf{4. Sanskritization fails to capture inverted reality} --- \textbf{Brahmins emulating Dalits}: \textbf{'Shudraization'} (a term by \textbf{Shyamlal}, a Rajasthan professor; upper castes wanting to become low caste for reservation benefits), which \textbf{D. N. Majumdar} calls \textbf{'desanskritization'}.
\end{itemize}
\end{CriticalPoint}

\section{Critique of Srinivas: Dominant Caste}

\begin{CriticalPoint}
\begin{itemize}
\item Two \textbf{strands of critique} on the role of numerical preponderance:
\item \textbf{Strand 1 --- numbers do not decide dominance}: \textbf{D. N. Majumdar} --- scheduled castes preponderate in many villages, yet the upper castes exercise power and authority (as headmen). \textbf{Adrian (A. C.) Mayer} coined \textbf{'village patriotism'} --- affiliation to the village is so strong that the numerical preponderance of the low castes is \textbf{subsumed to the development of the village}, even if led by upper castes (village solidarity overrides caste inequality; patriotism legitimises hierarchy).
\item \textbf{Strand 2 --- numbers do decide dominance}: \textbf{Beteille} --- before economic modernization the land-owning class (Brahmin) was powerful; with \textbf{land reforms} the big land-owning class \textbf{ceased to be} the source of dominance, and \textbf{numerical support became the decisive factor} (democracy is about numbers). Till 1946 a Brahmin ran the village panchayat; after \textbf{universal adult suffrage}, numbers mattered more than ritual status --- hence the \textbf{non-Brahmin movement} (and the rise of the \textbf{Kallar} sub-caste in Madurai politics; DMK/ADMK cabinets have 4--5 Kallar ministers).
\item In conclusion, \textbf{the criteria of dominance vary from village to village}.
\end{itemize}
\end{CriticalPoint}

\begin{CriticalPoint}
\textbf{Even the dominant castes now seek reservation} --- Marathas, Jats, Reddys, Patils, Patidars (feeling relative deprivation). '\textbf{They have to go down so that they can go up.}' This makes \textbf{desanskritization the new field view and Sanskritization the old (book) view} --- in the field we now see the dominant caste emulating the low caste (to get the backward tag). \textbf{Satish Deshpande} calls them '\textbf{backward forward}' (context makes them backward or forward).
\end{CriticalPoint}

\section{T. K. Oommen's Multiple Power Structures}

\begin{KeyConcept}
\begin{itemize}
\item \textbf{T. K. Oommen} (JNU Professor Emeritus): the concept of dominant caste \textbf{does not capture the alternate situations of dominance} --- e.g. a \textbf{numerically weak caste owning most of the land and wealth}, or a \textbf{numerically strong caste that is economically and ritually depressed}.
\item He therefore talks of \textbf{multiple power structures}: in a village there are \textbf{'power reservoirs'} (the pool of powerful people) and \textbf{'power exercisers'} (only the few owners of the reservoir actually exercise power).
\end{itemize}
\end{KeyConcept}

\section{Secularization of Caste and Caste Outside India}

\begin{KeyConcept}
\begin{itemize}
\item \textbf{Secularization of caste} (a concept of \textbf{D. L. Sheth}, alongside his 'classization of caste'): the \textbf{ritual function of caste is getting over}, and caste adopts a new \textbf{secular function of power} --- its modern avatar, the \textbf{vote bank}.
\item \textbf{Caste outside India / in every religion}: caste exists in every religion --- Christianity (Dalit Christians), Sikhism (Dalit Sikhs), and Islam (the \textbf{Arzal} = the lower/untouchable groups, vs the \textbf{Ashraf} = the upper groups).
\item \textbf{The function of caste in modern times is political} (the vote bank) --- earlier it had only a ritual function.
\end{itemize}
\end{KeyConcept}

\begin{QuickRevision}
\begin{itemize}
\item Ghurye = Orientalist/Indological/Evolutionary/Comparative (caste integrates); Srinivas = first official field view, structural-functional (Coorgs), Sanskritization (Varna) + dominant caste (Jati); Dumont = Indologist + structuralist, 'power inferior to status'; Beteille = critic (Brahmins not homogeneous).
\item Common to Ghurye/Srinivas/Dumont: Brahminical superiority; Beteille exposes it.
\item Critique of Sanskritization: legitimises Brahminical supremacy; ignores other traditions (Punjab got Islamized; Yogendra Singh; S. C. Dube's 'multiple traditions'); reference groups changed (secular/westernized; Dalit Assertion); inverted reality (Shudraization/Shyamlal; desanskritization/D. N. Majumdar).
\item Critique of dominant caste: numbers don't decide (D. N. Majumdar; Adrian Mayer's 'village patriotism') vs numbers do decide (Beteille, land reforms, universal suffrage, non-Brahmin movement, Kallar); criteria vary village to village.
\item Dominant castes seek reservation (Maratha/Jat/Reddy/Patil/Patidar) $\rightarrow$ desanskritization = new field view; Satish Deshpande's 'backward forward'.
\item T. K. Oommen: alternate situations of dominance; multiple power structures (power reservoir vs power exercisers).
\item Secularization of caste (D. L. Sheth) = ritual function over, secular function (vote bank); caste in every religion (Arzal/Ashraf; Dalit Christians/Sikhs).
\end{itemize}
\end{QuickRevision}

